\subsection{Adaptive Damage Localization via YOLOv8 Deep Learning}
\label{sec:crack_localization}

\begin{sloppypar}
The initial stage of our structural assessment pipeline involves the rapid localization of damage features within high-resolution imagery. We employ the state-of-the-art YOLOv8 (You Only Look Once) architecture, specifically the \texttt{yolov8n} (nano) variant, fine-tuned for the binary detection of masonry cracks. This stage serves as a critical spatial filter, isolating Regions of Interest (RoIs) for high-fidelity morphological analysis and reducing the computational overhead of downstream pixel-level processing.
\end{sloppypar}

\subsubsection{Architectural Paradigm: C2f and SPPF Modules}
\begin{sloppypar}
Unlike its predecessors, the YOLOv8 architecture utilizes a backbone characterized by the Cross-Stage Partial Bottleneck with two convolutions (\textbf{C2f}) module. This module enhances gradient flow and feature representation by concatenating the outputs of multiple bottleneck layers, allowing the network to capture complex crack textures at various scales. To provide a global receptive field, the architecture integrates a Spatial Pyramid Pooling - Fast (\textbf{SPPF}) layer at the bottleneck. The SPPF module performs successive max-pooling operations to extract multi-scale context, which is essential for distinguishing structural cracks from surface irregularities such as shadows or mortar joints.
\end{sloppypar}

\subsubsection{Training Protocol and Optimization}
\begin{sloppypar}
The model was optimized using the \textbf{Adam} optimizer with an initial learning rate $\eta = 0.001$, governed by a **Cosine Learning Rate Decay** schedule to ensure stable convergence. To enhance the model's robustness against the heterogeneous nature of disaster-site imagery, we implemented a heavy data augmentation pipeline, including:
\end{sloppypar}

\begin{itemize}[noitemsep]
    \item \textbf{Mosaic Augmentation:} Combining four training images into one to expose the model to varying object scales and occlusion patterns.
    \item \textbf{HSV Jittering:} Randomized adjustments to Hue, Saturation, and Value to simulate different lighting conditions.
    \item \textbf{Geometric Transforms:} Random rotations and flipping to ensure orientation-agnostic localization.
\end{itemize}

\begin{sloppypar}
The training objective minimizes a composite loss function $\mathcal{L}_{total}$, comprising Distribution Focal Loss ($\mathcal{L}_{dfl}$) for precise bounding box regression and Binary Cross-Entropy for classification.
\end{sloppypar}

\subsubsection{Performance Evaluation}
\noindent The localization model demonstrated exceptional performance on the validation set after 10 epochs of training. The specific metrics achieved are summarized below:

\begin{itemize}[noitemsep,leftmargin=1.5em]
    \item \textbf{Precision ($P$):} $0.933$
    \item \textbf{Recall ($R$):} $0.930$
    \item \textbf{Mean Average Precision ($mAP_{50}$):} $0.962$
    \item \textbf{Mean Average Precision ($mAP_{50-95}$):} $0.715$
\end{itemize}

\begin{sloppypar}
These high-fidelity results ensure that the pipeline can reliably anchor the diagnostic process on actual damage sites, effectively eliminating false positives from non-structural surface features.
\end{sloppypar}

\subsubsection{Dynamic Region of Interest (RoI) Generation}
\begin{sloppypar}
Each detected instance $d_j$ is represented by a bounding box $B_j = [x_1, y_1, x_2, y_2]$. To ensure sufficient context for the subsequent segmentation network, the localized boxes are dynamically expanded by a padding factor $\epsilon = 0.1 \times \max(w, h)$. This expansion prevents the "clipping" of crack tips, which are critical for determining propagation direction and structural risk. The resulting RoIs are cropped and normalized for the morphological extraction stage.
\end{sloppypar}
